\scanpage{102}
To tell the truth, it would not be difficult to require a condition 5°) which brings together 3°) and 4°), when one is given not $X'_1\to X'_2$, but
\diagram{\begin{tikzcd}[ampersand replacement=\&,row sep=3em,column sep=4em,
execute at end picture={
  \path (agA102.east) -- (agB102.west)
    coordinate[pos=.12] (agH102a) coordinate[pos=.35] (agH102b)
    coordinate[pos=.58] (agH102c) coordinate[pos=.82] (agH102d);
  \draw ([yshift=-.22em]agH102a) -- ([yshift=.32em]agH102b)
    -- ([yshift=-.22em]agH102b) -- ([yshift=.32em]agH102c)
    -- ([yshift=-.22em]agH102c) -- ([yshift=.32em]agH102d);
  \path (agB102.south) -- (agD102.north)
    coordinate[pos=.10] (agV102a) coordinate[pos=.32] (agV102b)
    coordinate[pos=.54] (agV102c) coordinate[pos=.76] (agV102d);
  \draw ([xshift=-.35em]agV102a) -- ([xshift=.35em]agV102a)
    -- ([xshift=-.35em]agV102b) -- ([xshift=.35em]agV102b)
    -- ([xshift=-.35em]agV102c) -- ([xshift=.35em]agV102c)
    -- ([xshift=-.35em]agV102d) -- ([xshift=.35em]agV102d);
}]
|[alias=agA102]| X'_1 \arrow[r,"\deleted{f_1}"] \arrow[dr,"f_1"] \&
|[alias=agB102]| \deleted{\unsure{Y_1}} \arrow[d,"\deleted{f}"] \\
X'_2 \arrow[r,"f_2"] \& |[alias=agD102]| Y_{\unsure{\deleted{2}|2}}
\end{tikzcd}}
with $f_2$ \deleted{\unsure{étale}} inducing a homomorphism \deleted{\unsure{fibrant}} surjective on the $\pi_1$, so that $\widehat{\Pi}_1(Y,f_2(I'_2))$ can be interpreted as a quotient of $\widehat{\Pi}_1(X'_2,I'_2)$, so that the action \deleted{of} $\Theta'_2$ of $\Theta$ on the latter descends to $\widehat{\Pi}_1(Y,f_2(I'_2))$ \ldots

But I am not sure whether I shall need this common generalization of 3° and 4° (and $=$ of 2°, 3°, 4°, to tell the truth!).

Conditions 2° to 4° are for me above all a convenient \uline{recipe} for finding necessary conditions on the virtual parameters which specify an undetermined $u\in\Gamma$, for them to be effective i.e. really do come from a $u$. The point of view \inserted{here} is not at all to bring out \deleted{some} necessary and sufficient \unsure{condition|conditions} which will make it possible to \uline{characterize} the $\Theta$ coming from $\Gamma$ \ldots
\mspage{101}
